\section{Limitations and Future Work}\label{sec:limitations}

\subsection{Limitations}

The Bailout Protocol provides a structurally compulsive escalation pathway from any \bxthree{} node to a named human principal, but four structural constraints limit its operation in adversarial, high-latency, or large-scale deployment contexts. We state each limitation precisely, identify its operative scope, and describe the mitigation pathway where one exists.

\subsubsection{Human Response Latency}

The Bailout Protocol's liveness property holds under the assumption that the human anchor $h(n)$ is reachable and responsive within the bound $T_{\text{human}}$. In practice, $h(n)$ may be unavailable --- engaged in other tasks, offline, or incapacitated --- for periods that exceed the configured timeout. When $T_{\text{human}}$ expires without a resolution directive, the protocol's current specification transitions to \texttt{RESOLVED}$(h(n), \bot)$, recording a null resolution with a protocol violation flag.

This outcome is architecturally correct under the timeout model but practically insufficient for systems requiring continuous operation without human presence. A null resolution with a violation flag surfaces the failure to auditors but does not restore the node to productive operation. The node remains suspended in \texttt{ESCALATING} until $h(n)$ responds, creating a denial-of-service condition on the node's execution that is produced by the protocol itself rather than by the original exception. For systems requiring uninterrupted operation --- critical infrastructure, autonomous vehicles, real-time financial systems --- this limitation implies that the Bailout Protocol must be deployed alongside a secondary human-anchor pool or a graceful degradation pathway that the current specification does not provide.

The latency limitation is not correctable through timeout reduction alone. Decreasing $T_{\text{human}}$ increases the probability of premature termination (capturing a responsive human as unavailable), while increasing it increases the vulnerable window during which a suspended node cannot contribute to the system's operation. The optimal timeout value is both context-dependent and dynamically varying --- dependent on $h(n)$'s current workload and availability, and varying as operational conditions change.

\subsubsection{Adversarial Conditions}

The Bailout Protocol assumes that the exception conditions triggering escalation are authentic --- that a node has genuinely encountered a state outside its operating range $R(n)$ that it cannot resolve within the bounded reasoner. Under adversarial conditions, this assumption does not hold. A compromised node, or a node operating under a adversarial purpose-layer injection, can fabricate exception conditions deliberately to trigger the Bailout Protocol for strategic purposes: to exhaust human attention, to trigger denial-of-service on specific nodes, or to map the escalation chain for subsequent targeted interference.

The protocol's bypass mechanism propagates the exception upward regardless of whether the exception is authentic, because the \bounds{} layer cannot distinguish authentic from fabricated boundary violations. The Forensic Ledger records each escalation event with full attribution, but attribution to a named node does not establish that the node's reported exception reflects the node's actual operational state. An adversarial actor with control over a node's \bounds{} layer can produce the same observable escalation traces as an authentic boundary violation, and the protocol provides no mechanism to detect the difference at escalation time.

This limitation is particularly acute in shared or multi-tenant \bxthree{} deployments where not all nodes are operated by principals with aligned interests. Without an attestation mechanism that independently verifies the node's operational state at the \fact{} layer, the Bailout Protocol's escalation pathway can be weaponized by any actor with node-level control. The protocol's structural compulsion --- its defining strength in non-adversarial contexts --- becomes a surface for adversarial exploitation under hostile deployment conditions.

\subsubsection{Scalability of the Human Anchor}

The Bailout Protocol establishes a one-to-one mapping between each node and its human anchor $h(n)$. In a \bxthree{} system with $N$ active nodes, this implies $N$ independent escalation endpoints, each requiring a distinct named human principal with the authority and availability to issue binding resolution directives. For small-to-medium agent populations ($N < 50$), this mapping is operationally tractable: a single human principal can serve as $h(n)$ for multiple nodes within a well-defined domain, and the escalation load remains within the bounds of human cognitive throughput.

For large agent populations ($N \gg 100$), the one-to-one mapping becomes a structural bottleneck. If every node in a population of 1,000 agents can independently escalate to the same human anchor, that anchor faces a potential influx of 1,000 simultaneous or near-simultaneous escalation events under a system-wide exception condition. The current protocol provides no mechanism for aggregating multiple simultaneous escalations into a single human-facing event, nor does it provide any load-balancing or priority-ranking mechanism across concurrent escalation requests addressed to the same $h(n)$. All concurrent escalations to $h(n)$ arrive as independent, unranked events, imposing a coordination problem on the human anchor that the protocol's design does not address.

The scalability limitation is not merely theoretical. Real-world autonomous systems operate populations of hundreds to thousands of concurrent agents under a single organizational authority. A \bxthree{} deployment at this scale without a corresponding scaling of the human anchor infrastructure would produce an escalation load that exceeds human cognitive throughput, defeating the protocol's liveness property by overwhelming the human anchor rather than by timing out.

\subsubsection{Exploitation of the Escalation Pathway}

Closely related to the adversarial conditions limitation but distinct in its mechanism is the risk of strategic exploitation of the escalation pathway itself. Even in the absence of node-level compromise, an adversarial actor with knowledge of the \bxthree{} architecture and the Bailout Protocol's specification can craft inputs designed to trigger repeated, low-severity exceptions that produce continuous escalation without satisfying any individual severity threshold that might warrant system-level intervention. The protocol's sensitivity to observational boundary violations --- its trigger mechanism --- is calibrated for unanticipated exceptions, not for deliberately crafted inputs that exploit the gap between boundary sensitivity and severity.

This exploitation pathway resembles a slow-dose attack: each individual exception is within the human anchor's capacity to resolve, but the accumulated cognitive load of resolving a high frequency of low-severity escalations degrades the human anchor's ability to respond to genuinely critical escalations when they occur. The protocol provides no mechanism for rate-limiting escalation frequency, aggregating repeated escalations from the same node, or distinguishing escalation patterns that indicate systemic malfunction from those that indicate deliberate exploitation.

The Forensic Ledger records each escalation with full attribution, enabling post-hoc pattern analysis that might detect exploitation, but the protocol itself provides no runtime mechanism to act on this detection. The human anchor must infer the pattern from the ledger and issue a corrective directive, but the protocol offers no structured way to prioritize, suppress, or respond to a detected exploitation pattern without a dedicated administrative action.

\subsection{Future Work}

Each of the four limitations identified above suggests a directed research agenda. We describe three priority directions that address the most operationally significant gaps in the protocol's current specification.

\subsubsection{Formal Verification of the Protocol}

The Bailout Protocol's correctness properties --- Safety, Liveness, and Accountability --- are currently stated as theorems with proof sketches derived from the formal state machine specification. A more rigorous treatment requires machine-checked formal verification using a verification-aware specification system such as TLA+\cite{lamport2019tla}, Ivy\cite{ivy}, or Coq\cite{coq}. Formal verification would establish the correctness properties under the stated assumptions with mathematical certainty, eliminate the possibility of logical gaps in the proof sketches, and provide a basis for verifying concrete implementations against the verified specification.

The verification agenda has two components. The first is the verification of the protocol's state machine in isolation: that the transition relation is total over the defined state-event pairs, that every terminal state satisfies either \texttt{RESOLVED} or \texttt{IDLE}, and that the bypass mechanism is preserved under all reachable states. The second is the verification of the protocol's embedding within the \bxthree{} three-layer architecture: that the \fact{} layer's enforcement of the transition relation is consistent with the verified state machine, that the immutable parent pointer cannot be subverted by any operation within the \bounds{} or \purpose{} layers, and that the Forensic Ledger's append-only property is preserved under concurrent access.

Formal verification is a prerequisite for deployment of the Bailout Protocol in safety-critical or regulatory-mandated contexts, where the correctness of the escalation mechanism must be demonstrable to external auditors or regulatory bodies rather than merely asserted by the system designer.

\subsubsection{Adaptive Timeout and Escalation Aggregation}

The human response latency and scalability limitations both point toward the need for an adaptive timeout mechanism that adjusts $T_{\text{human}}$ dynamically based on the current state of the human anchor pool and the aggregate escalation load across the agent population. Rather than a fixed timeout per escalation event, an adaptive mechanism would monitor the human anchor's current response latency, the number of unresolved escalations currently pending for each human anchor, and the estimated severity of each pending escalation, and would adjust timeout values and escalation aggregation accordingly.

An adaptive timeout mechanism requires three supporting capabilities that the current protocol does not specify. First, a monitoring subsystem that tracks $h(n)$'s response latency across recent escalation events and maintains a rolling estimate of the anchor's current availability. Second, an aggregation subsystem that groups concurrent escalations from the same or related nodes into a single human-facing event with a unified severity assessment, reducing the cognitive load imposed by simultaneous multi-node failures. Third, a priority subsystem that ranks concurrent escalations by severity and relevance to the human anchor's domain, ensuring that the most critical escalations receive human attention first regardless of arrival order.

The design of an adaptive mechanism must preserve the protocol's non-discretionary entry property: escalation remains compulsory and cannot be suppressed or deferred by the node or by any machine actor. Adaptation applies to the timeout value and to the aggregation structure, not to the obligation to escalate.

\subsubsection{Distributed Human Anchor Pools}

The scalability limitation is most directly addressed by replacing the one-to-one node-to-anchor mapping with a distributed human anchor pool: a set of two or more human principals who share escalation responsibility for a node or a domain of nodes, with a consensus or load-balancing mechanism for routing escalations to available anchors. A distributed anchor pool eliminates the single-point-of-failure produced by the one-to-one mapping, allows the human anchor capacity to scale with the agent population, and provides resilience against individual anchor unavailability.

A distributed anchor pool introduces its own correctness requirements that the current protocol does not address. The pool must present to the protocol as a single terminus --- the bypass mechanism's requirement that no machine actor appear between the node and the human anchor must be preserved even when multiple human principals share anchor responsibility. This requires a pool coordination layer that is itself architecturally constrained: the coordination mechanism must not be a machine actor capable of absorbing exceptions, and the consensus or routing decision within the pool must not be able to suppress or redirect an escalation against the pool's configured policy.

A distributed anchor pool also introduces accountability complexity: when an escalation is resolved by a specific human principal within the pool, the resolution must be attributed to that individual with the same precision as a single-anchor resolution, and the pool's internal routing decision must be recorded in the Forensic Ledger. The Accountability property of the Bailout Protocol --- that $h(n)$ is always identifiable at the terminus --- extends naturally to pool configurations, but requires a formal treatment of pool membership, pool routing policy, and individual attribution that the current specification does not provide.

\subsection{Research Roadmap}

The three future work directions form a research roadmap ordered by dependency. Formal verification (Direction 1) is a prerequisite for deployment in safety-critical contexts and provides the specification certainty required to design adaptive mechanisms rigorously. Adaptive timeout and escalation aggregation (Direction 2) addresses the latency and scalability limitations for single-anchor deployments and provides the monitoring and aggregation infrastructure needed for pool designs. Distributed human anchor pools (Direction 3) provides the long-term scalability path for large agent populations and multi-tenant deployments. Together, these directions extend the Bailout Protocol from a structurally sound specification toward an operationally robust escalation system suitable for real-world deployment at scale.
